\section{Peeling the live core while preserving physical-flux groups}
\label{sec:op3-live-core-peeling}

\paragraph{Status and scope.}
The following \statusproved{} construction is a proof draft awaiting
independent review. It strengthens \cref{thm:op3-branch-core-flux} by
allowing retained vertices to disappear after their side branches settle.
The work depends on the maximum simultaneous core size, including a
temporarily admitted vertex, rather than on the total number of original
branching vertices admitted. No graph-uniform bound on that maximum is
claimed. General OP3 remains \statusopen.

Use $\lambda=\epsappr/2$, $\bm M=\bm D-\gamma\bm A$, and
$\bm b=\bm e_v-\lambda\bm d$. Every strictly positive admission is
temporarily retained. The seed always remains retained. Other core vertices
are eliminated whenever their current reduced degree is at most two,
counting distinct retained and inactive neighbors after parallel couplings
have been coalesced. Denote the current retained set by $C$ and exposed
inactive set by $B$. The admitted original set $U$ includes both retained
and eliminated vertices.

\subsection{The reduced matrix and grouped physical signals}

Maintain the exact Schur matrix on $C\cup B$, retaining its core block,
core--boundary couplings, boundary diagonals and transformed loads.
Original inactive--inactive edges are revealed only on admission of an
endpoint. No fill edge between two inactive vertices is created.
For a pair $(c,j)\in C\times B$, collect a disjoint group of original
boundary edges whose total physical flux has the form
\begin{equation}
 F_{cj}=a_{cj}u_c+b_{cj},\qquad
 a_{cj}>0,\quad b_{cj}\leq0.
 \label{eq:op3-grouped-physical-flux}
\end{equation}
Store the original incidence count $k_{cj}$ and a lower publication
$\ell_{cj}$. At each inactive $j$ the groups partition all original edges
from $U$ into $j$, so its exact gate is
\[
 g_j=-\lambda d_j+\sum_c F_{cj}.
\]
The group slope is also the reduced coupling between $c$ and $j$.
Its intercept contributes to the transformed load of $j$.

\begin{lemma}[One boundary group per eligible core removal]
\label{lem:op3-live-core-removal}
If a nonseed core vertex $i$ has reduced degree at most two, it has at
least one core neighbor and at most one inactive neighbor. Eliminating
$i$ changes at most two remaining neighbors and moves at most one
physical-flux group. The representation
\cref{eq:op3-grouped-physical-flux}, its signs, and the absence of
inactive--inactive fill are preserved.
\end{lemma}
\begin{proof}
Every admitted vertex has an original active neighbor at admission, hence
the admitted graph is connected to the seed. Positive Schur elimination
preserves connectedness among the remaining retained vertices. Thus a
nonseed retained vertex has a retained neighbor, proving the degree claim.

Suppose the two neighbors are retained $c$ and inactive $j$. Write
\[
 \delta_i u_i=b_i+w_{ic}u_c+a_{ij}u_j,
 \qquad F_{ij}=a_{ij}u_i+b_{ij}.
\]
At the pinned boundary value $u_j=0$, elimination moves this physical group
to retained coordinate $c$ with
\begin{equation}
 a'_{cj}=a_{ij}w_{ic}/\delta_i,\qquad
 b'_{cj}=b_{ij}+a_{ij}b_i/\delta_i,
 \qquad \delta_j\leftarrow\delta_j-a_{ij}^2/\delta_i.
 \label{eq:op3-live-core-group-move}
\end{equation}
All nonseed transformed core loads remain negative: initially they are
$-\lambda d_i$, and eliminating another negative-load vertex adds a
negative multiple to them. Schur pivots remain positive. Thus the group
signs are preserved. Keep its original incidence count and lower
publication unchanged. If a group $(c,j)$ already exists, add the slopes,
intercepts, incidence counts and lower publications. The physical edge
sets being combined are disjoint.

For two retained neighbors, ordinary two-neighbor Schur complementation
changes their diagonals, loads and mutual coupling; there is no boundary
group to move. With one retained neighbor and no inactive neighbor, the
operation removes a settled leaf response. In all cases there are at most
two distinct neighbors, so no two inactive vertices receive a fill edge.
Save the removed scalar equation for recovery.
\end{proof}

When $j\in B$ is admitted, its current diagonal, transformed load
$-\lambda d_j+\sum_c b_{cj}$, and couplings $a_{cj}$ form its exact
bordered core equation. Retire all groups ending at $j$, then scan its
original row. A newly revealed edge to an inactive vertex starts a direct
group with $(a,b,k,\ell)=(\gamma,0,1,0)$. An original-edge dictionary
prevents repeating an old incidence. Queue the new core vertex for peeling;
after a removal, only its retained neighbors need another eligibility test.
Thus no global scan is needed to find eligible removals.

\subsection{Events, inverse restriction and recovery}

Use a scalar minimum heap at each retained coordinate, with group threshold
\begin{equation}
 \tau_{cj}=\frac{(5/4)\ell_{cj}+k_{cj}\epsappr/16-b_{cj}}{a_{cj}}.
 \label{eq:op3-grouped-flux-threshold}
\end{equation}
After updating the core, settle all these heaps. If $u_c>\tau_{cj}$,
publish the actual group flux and update the shared sum
$L_j=\sum_c\ell_{cj}$. Queue $j$ once this sum exceeds
$(11/20)\epsappr d_j$. Equality stays inactive.

Moving and coalescing groups preserves the published sum at their common
target. In particular, a moved group's lower publication is not reset to
zero. A published group remains a lower bound after every legal admission
by monotonicity of the restricted solution. A settled group satisfies
\[
 \ell_{cj}\leq F_{cj}\leq\tfrac54\ell_{cj}
                  +k_{cj}\epsappr/16.
\]
Summing over the disjoint physical incidences gives the same original
residual certificate as \cref{eq:op3-arm-terminal-residual}. Every queued
candidate stays legal, including through representation changes.

Maintain the retained inverse $\bm P=\bm H_C^{-1}$ and retained values.
An admission uses the bordered update in
\cref{eq:op3-core-rank-update}. Removing a retained vertex leaves the
restricted solution unchanged, and the inverse of its new Schur matrix is
exactly the principal submatrix of $\bm P$ on the remaining retained labels.
The implementation stores dense inverse rows keyed by original vertex
labels. Deleting one row and column costs $O(|C|)$ dictionary entries;
there is no global change of the labels used by old group or recovery records.

Each original positive vertex is either still retained or has exactly one
saved equation. Its coefficients name original remaining-neighbor labels,
including any inactive neighbor at that time. In reverse elimination order,
every later eliminated neighbor has already been recovered, and a neighbor
that never activated contributes zero. This recovers all original positive
coordinates once. A vertex's temporary core birth is not a second output
record.

\begin{theorem}[A bound in maximum live-core size]
\label{thm:op3-live-core-peeling}
Let $q$ be the largest simultaneous retained set, including temporary
admissions before peeling. Put $V=\vol(U)$, $n_U=|U|$, and
$L=1+\log_+(1/(\bar\alpha\epsappr))$ for the returned support $U$.
On every canonical input graph, the algorithm returns a nonnegative ACL
approximation with exact-real word work
\begin{equation}
 O\left((n_Uq^2+qV+VL)\log(2+V)\right)
 =\widetilde O((1+q^2)/\epsappr),
 \label{eq:op3-live-core-work}
\end{equation}
and working storage $O(V+q^2)$. Empty output costs $O(1)$.
Only positive-output adjacency rows are scanned, and degrees are queried
only at the seed or exposed neighbors. No topology, elimination order,
final support or bound on $q$ is supplied.
\end{theorem}
\begin{proof}
The reduced equations and \cref{lem:op3-live-core-removal} preserve the exact
restricted solution. Scalar publications certify positive gates for the
obstacle at $\lambda=\epsappr/2$. Inverse positivity therefore makes all
admissions safe and all old coordinates increase. At termination active
residuals equal $\lambda d_i$, while the summed group bands bound inactive
residuals by $(3/4)\epsappr d_i$. This yields the ACL witness, and
conservation gives $V\leq2/\epsappr$.

Each original row is scanned once. Each admitted vertex is removed from
the core at most once. A removal changes only two reduced neighbors and
moves at most one group. A direct edge starts one group and each move
starts at most one replacement group, possibly coalescing an existing
group. Thus the total number of group lifetimes is $O(V)$, including all
replacements. During a fixed lifetime, physical membership is fixed and
$0\leq F_{cj}\leq k_{cj}\gamma/\bar\alpha$. A first new publication
exceeds $k_{cj}\epsappr/16$, and successive publications grow by more
than $5/4$. Hence each lifetime has $O(L)$ publications, regardless of its
incidence count. Each retired entry is discarded once, including heaps
discarded with a removed core. Every unsuccessful heap test, core-queue
scan, ready-queue operation and peeling eligibility test is charged.

The number of core-queue visits is at most $n_Uq$. Each admission updates
at most $q^2$ inverse entries; sparse inverse-vector products cost at most
$O(qV)$ in total because candidate group counts sum to $O(V)$. Removing
inverse rows and columns costs $O(n_Uq)$ entries. Reduced-edge changes,
original-edge bookkeeping, group replacement and reverse equations cost
$O(V)$ records. Comparison dictionaries and scalar heaps give the displayed
logarithmic factor and a deterministic word bound. The audited Python code
uses dictionaries as a bookkeeping convenience. Dense inverse storage is
$O(q^2)$, and all group lifetimes, cached rows, scalar heaps and recovery
records require $O(V)$ words. No group operation copies its original
physical membership list; such lists appear only in the independent audit.
\end{proof}

\subsection{A source-valid separation between two admission orders}

\begin{proposition}[Peeling resolves the balanced-tree diagnostic under depth-first order]
\label{prop:op3-live-core-binary-order}
On the complete binary tree of height $h\geq2$, rooted at the seed, take
$\alpha=1/1009$ and $\epsappr=\alpha/(10\cdot6^h)$ as in
\cref{prop:op3-dense-core-binary-tree}. The last-in-first-out ready policy
has $q\leq h+1$ and therefore $\widetilde O(n)$ word work for this family.
The first-in-first-out policy has $q\geq2^h$ and incurs
$\Omega(n^3)$ explicit inverse writes before the first leaf is processed.
Both policies return valid ACL output on the same full support.
\end{proposition}
\begin{proof}
Every newly exposed boundary vertex becomes ready after the heaps settle.
To see this uniformly over connected root-containing faces, let $i$ be
the active parent of an inactive boundary vertex. Its depth is at most
$h-1$. A single walk term gives
$(\bm M_{UU}^{-1}\bm e_v)_i\geq2/6^h$.
Also $\bm M_{UU}\bm1\geq\bar\alpha\bm d_U$, so
$\bm M_{UU}^{-1}\bm d_U\leq\bm1/\bar\alpha$. Consequently
\[
 u_i^U\geq\frac2{6^h}-\frac\lambda{\bar\alpha}
 >\frac{39}{20\cdot6^h},\qquad
 r_j^U\geq\gamma u_i^U>\frac{39}{40\cdot6^h}
 >\tfrac34\epsappr d_j.
\]
The settled group bands then force
$L_j>(11/20)\epsappr d_j$. All old boundary vertices were already queued;
the newly queued vertices are exactly the children exposed by the latest
admission. Thus the two ready policies realize the usual breadth-first
and depth-first orders, with arbitrary sibling ties.

During depth-first exploration of a tree, all active vertices off the
current root path belong to completed subtrees. Their leaf responses have
been peeled, recursively, and their roots are eligible for further peeling.
No unfinished active subtree lies off the current search path. All retained
vertices therefore lie on that path; even including a temporary admission
there are at most $h+1$. Since $V=O(n)$ and $h=O(\log n)$, the work bound
follows from \cref{thm:op3-live-core-peeling}.

Under breadth-first order every internal vertex is admitted before any
leaf. Each nonseed internal vertex still has three distinct reduced
neighbors, so none has been peeled. The first leaf is temporarily added to
the $2^h-1$ internal retained vertices, proving the bound on $q$. The
preceding internal births alone force
$\sum_{k=1}^{2^h-2}k^2=\Omega(n^3)$ inverse writes. This is a difference
between two implementations on an easy graph family, not an OP3 lower bound.
\end{proof}

\paragraph{What remains open.}
Depth-first order alone does not bound $q$ on a general trace. A branch that
is quiet at one root value can become productive after another branch
changes the response. Many unfinished branches may therefore remain live.
The next target is a paid rule that either bounds the simultaneous core,
eliminates a higher-degree core vertex with a stronger grouped response,
or maintains its response and all due thresholds implicitly. Every change
of physical membership and every failed certificate must remain charged.

\begin{proposition}[Interior tree support blocks every first peeling step]
\label{prop:op3-live-core-interior-tree}
For any integer $h\geq1$, let the input be the complete binary tree of
height $H=3h+5$, seeded at its root, with
\[
 \alpha=\tfrac13,\qquad \gamma=\bar\alpha=\tfrac12,
 \qquad \epsappr=\frac1{12\cdot6^h},\qquad
 \lambda=\epsappr/2.
\]
Under either ready policy, and indeed under any legal ready policy in this
construction, no nonseed core vertex is peeled. Every valid ACL output
contains all vertices through depth $h$, so
$q=|U|\geq2^{h+1}-1$. The explicit inverse backend consequently requires
$\Omega(8^h)=\Omega((1/\epsappr)^{\log_6 8})$ word writes.
In particular, current-degree-two peeling plus depth-first order does not
give a uniform polylogarithmic live core or an OP3 work bound.
\end{proposition}
\begin{proof}
A source-to-vertex walk term, using original degrees at most three, gives
for every vertex at depth at most $h$
\[
 \pi_i/d_i=\bar\alpha(\bm M^{-1})_{iv}
 \geq\frac1{6\cdot6^h}=2\epsappr.
\]
The ACL semantic implication therefore forces positive output at every
such vertex.

At the other end of the tree, no original leaf can belong to the positive
obstacle support at $\lambda$. To verify this, write
\[
 \bm M^{-1}=(\bm I-\gamma\bm D^{-1}\bm A)^{-1}\bm D^{-1}.
\]
The transition matrix is row-stochastic. Its entries connecting a vertex
at depth $t$ to the seed vanish before time $t$ and are at most one
afterward. Since $d_v\bar\alpha=1$, the unregularized degree potential
at that vertex is at most $\gamma^t$. The obstacle potential is bounded
above by this unregularized solution. At the parent of an original leaf,
the possible incoming leaf flux is thus at most
\[
 \gamma^H=\frac1{32\cdot8^h}
 <\frac1{24\cdot6^h}=\lambda.
\]
A positive leaf would require $u_i=\gamma u_{p(i)}-\lambda>0$, a
contradiction. Safe admissions consequently never reach any original leaf.

Before any peeling step, every admitted nonseed vertex has its three
distinct original neighbors in the core or the exposed boundary. An
admission changes a neighbor's status but does not reduce that degree.
There can therefore be no first eligible removal. Induction gives no
removals at all, and $C=U$. Each successive admission changes every entry
of the previous positive retained inverse, giving
$\sum_{k=1}^{|U|-1}k^2=\Omega(|U|^3)$. The required depth-$h$ prefix
proves the stated bound. Since $\log_6 8>1$ and all relevant logarithms
grow only with $h$, this exceeds the OP3 scale for this backend.
\end{proof}

The last proposition is an implementation-specific negative control.
Teleportation is fixed in that family, so standard positive local-push
bounds already reach $O(1/\epsappr)$ there. It neither refutes OP3 nor
rules out a method choosing between response representations. It does show
that a small live-core assertion requires additional structure, and that
completed-subtree examples alone are insufficient evidence for it.
